\documentclass[12pt,letterpaper]{hmcpset}
\usepackage{lin-alg-preamble}

% info for header block in upper right hand corner
\name{Jayden Thadani}
\class{Linear algebra done right}
\assignment{Exercises 6.A --- Inner Products and Norms}

\begin{document}

\begin{problem}[1]
	Show that the function that takes $((x_1, x_2), (y_1, y_2)) \in \bb{R}^2 \times \bb{R}^2$ to $\abs{x_1 y_1} + \abs{x_2 y_2}$ is not an inner product on $\bb{R}^2$.
\end{problem}

\pagebreak

\begin{problem}[2]
	Show that the function that takes $((x_1, x_2, x_3), (y_1, y_2, y_3)) \in \bb{R}^3 \times \bb{R}^3$ to $x_1 y_1 + x_3 y_3$ is not an inner product on $\bb{R}^3$. 
\end{problem}

\pagebreak

\begin{problem}[3]
	Suppose $\bb{F} = \bb{R}$ and $V \neq \set{0}$. Replace the positivity condition (which states that $\langle v, v \rangle \ge 0$ for all $v \in V$) in the definition of an inner product with the condition that $\langle v, v \rangle > 0$ for some $v \in V$. Show that this change in the definition does not change the set of functions from $V \times V$ to $\bb{R}$ that are inner products on $V$.
\end{problem}

\pagebreak

\begin{problem}[4]
	Suppose $V$ is a real inner product space.
	\begin{enumerate}
		\item Show that $\langle u + v, u - v \rangle = \norm{u}^2 - \norm{v}^2$ for every $u, v \in V$.
		\item Show that if $u, v \in V$ have the same norm, then $u + v$ is orthogonal to $u - v$.
		\item Use part (b) to show that the diagonals of a rhombus are perpendicular to each other.
	\end{enumerate}
\end{problem}

\pagebreak

\begin{problem}[5]
	Suppose $T \in \mathcal{L}(V)$ (for finite-dimensional $V$) is such that $\norm{Tv} \le \norm{v}$ for every $v \in V$. Prove that $T - \sqrt{2} I$ is invertible.
\end{problem}

\pagebreak

\begin{problem}[6]
	Suppose $u, v \in V$. Prove that $\langle u, v \rangle = 0$ if and only if
	\[
		\norm{u} \le \norm{u + av}
	\]
	for all $a \in \bb{F}$.
\end{problem}

\pagebreak

\begin{problem}[7]
	Suppose $u, v \in V$. Prove that $\norm{au + bv} = \norm{bu + av}$ for all $a, b \in \bb{R}$ if and only if $\norm{u} = \norm{v}$.
\end{problem}

\pagebreak

\begin{problem}[8]
	Suppose $u, v \in V$ and $\norm{u} = \norm{v} = 1$ and $\langle u, v \rangle = 1$. Prove that $u = v$.
\end{problem}

\pagebreak

\begin{problem}[9]
	Suppose $u, v \in V$ and $\norm{u} \le 1$ and $\norm{v} \le 1$. Prove that
	\[
		\sqrt{1 - \norm{u}^2} \sqrt{1 - \norm{v}^2} \le 1 - \abs{\langle u, v \rangle}
	\]
\end{problem}

\pagebreak

\begin{problem}[10]
	Find vectors $u, v \in \bb{R}^2$ such that $u$ is a scalar multiple of $(1, 3)$, $v$ is orthogonal to $(1, 3)$, and $(1, 2) = u + v.$
\end{problem}

\pagebreak

\begin{problem}[11]
	Prove that
	\[
		16 \le (a + b + c + d)(\frac{1}{a} + \frac{1}{b} + \frac{1}{c} + \frac{1}{d})
	\]
	for all positive numbers $a, b, c, d$.
\end{problem}

\pagebreak

\begin{problem}[12]
	Prove that
	\[
		(x_1 + \cdots + x_n)^2 \le n (x_1^2 + \cdots + x_n^2)
	\]
	for all positive integers $n$ and all real numbers $x_1, \ldots, x_n$.
\end{problem}

\pagebreak

\begin{problem}[13]
	Suppose $u, v$ are nonzero vectors in $\bb{R}^2$. Prove that
	\[
		\langle u, v \rangle = \norm{u} \norm{v} \cos{\theta}
	\]
	where $\theta$ is the angle between $u$ and $v$ (thinking of $u$ and $v$ as arrows with initial point at the origin).
\end{problem}

\pagebreak

\begin{problem}[14]
	The angle between two vectors (thought of as arrows with initial point at the origin) in $\bb{R}^2$ or $\bb{R}^3$ can be defined geometrically. However, geometry is not as clear for $\bb{R}^n$ for $n < 3$. Thus, the angle between two nonzero vectors $x, y \in \bb{R^n}$ is defined to be
	\[
		\arccos{\frac{\langle x, y \rangle}{\norm{x} \norm{y}}}
	\]
	where the motivation for this definition comes from the previous exercise. Explain why the Cauchy-Schwartz inequality is needed to show that this definition makes sense.
\end{problem}

\pagebreak

\begin{problem}[15]
	Prove that
	\[
		\Big(\sum_{j = 1}^n a_j b_j \Big) \le \Big(\sum_{j = 1}^n j a_j^2 \Big) + \Big(\sum_{j = 1}^n \frac{b_j^2}{j} \Big) 
	\]
	for all real numbers $a_1, \ldots, a_n$ and $b_1, \ldots, b_n$.
\end{problem}

\pagebreak

\begin{problem}[16]
	Suppose $u, v \in V$ are such that
	\[
		\begin{gathered}
			\norm{u} = 3 \\
			\norm{u + v} = 4 \\
			\norm{u - v} = 6.
		\end{gathered}
	\]
	What number does $\norm{v}$ equal?
\end{problem}

\pagebreak

\begin{problem}[17]
	Prove or disprove: there is an inner product on $\bb{R}^2$ such that the associated norm is given by
	\[
		\norm{(x, y)} = \max \set{\abs{x}, \abs{y}}
	\]
	for all $(x, y) \in  \bb{R}^2$
\end{problem}

\pagebreak

\begin{problem}[18]
	Suppose $p > 0$. Prove that there is an inner product on $\bb{R}^2$ such that the associated norm is given by
	\[
		\norm{(x, y)} = (\abs{x}^p + \abs{y}^p)^{1/p}
	\]
	for all $(x, y) \in \bb{R}^2$ if and only if $p = 2$.
\end{problem}

\pagebreak

\begin{problem}[19]
	Suppose $V$ is a real inner product space. Prove that
	\[
		\langle u, v \rangle = \frac{\norm{u + v}^2 - \norm{u - v}^2}{4}	
	\]
	for all $u, v \in V$.
\end{problem}

\pagebreak

\begin{problem}[20]
	Suppose $V$ is a complex inner product space. Prove that
	\[
		\langle u, v \rangle = \frac{\norm{u + v}^2 - \norm{u - v}^2 + \norm{u + i v}^2 i - \norm{u - i v}^2 i}{4}	
	\]
	for all $u, v \in V$.
\end{problem}

\pagebreak

\begin{problem}[21]
	A norm on a vector space $U$ is a function $\norm{\,}: U \to [0, \infty)$ such that $\norm{u} = 0$ if and only if $u = 0$, $\norm{\alpha u} = \abs{\alpha} \norm{u}$ for all $\alpha \in \bb{F}$ and all $u \in U$, and $\norm{u + v} \le \norm{u} + \norm{v}$ for all $u, v \in U$. \\\\
	Prove that a norm satisfying the parallelogram equality comes from an inner product (in other words, show that if $\norm{\,}$ is a norm on $U$ satisfying the parallelogram equality, then there is an inner product $\langle \,\, , \, \rangle$ on $U$ such that $\norm{u} = \langle u, u \rangle^{1/2}$ for all $u \in U$).
\end{problem}

\pagebreak

\begin{problem}[22]
	Show that the square of an average is less than or equal to the average of the squares. More precisely, show that if $a_1, \ldots, a_n \in \bb{R}$, then the square of the average of $a_1, \ldots, a_n$ is less than or equal to the average of $a_1^2, \ldots, a_n^2$.
\end{problem}

\pagebreak

\begin{problem}[23]
	Suppose $V_1, \ldots, V_m$ are inner product spaces. Show that the equation
	\[
		\langle (u_1, \ldots, u_m), (v_1, \ldots, v_m) \rangle = \langle u_1, v_1 \rangle + \cdots + \langle u_m, v_m \rangle
	\]
	defines an inner product on $V_1 \times \cdots \times V_m$
\end{problem}

\pagebreak

\begin{problem}[24]
	Suppose $S \in \mathcal{L}(V)$ is an injective operator on $V$. Define $\langle \cdot, \cdot \rangle_1$ by
	\[
		\langle u, v \rangle_1 = \langle Su, Sv \rangle
	\]
	for $u, v \in V$. Show that $\langle \cdot, \cdot \rangle_1$ is an inner product on $V$.
\end{problem}

\pagebreak

\begin{problem}[25]
	Suppose $S \in \mathcal{L}(V)$ is not injective. Define $\langle \cdot, \cdot \rangle_1$ as in the exercise above. Explain why $\langle \cdot, \cdot \rangle_1$ is not an inner product on $V$.
\end{problem}

\pagebreak

\begin{problem}[26]
	Suppose $f, g$ are differentiable functions from $\bb{R}$ to $\bb{R}^n$.
	\begin{enumerate}
		\item Show that
			\[
				\langle f(t), g(t) \rangle' = \langle f'(t) g(t) \rangle + \langle f(t), g'(t) \rangle.
			\]
		\item Suppose $c > 0$ and $\norm{f(t)} = c$ for every $t \in \bb{R}$. Show that $\langle f'(t), f(t) \rangle = 0$ for every $t \in \bb{R}$.
		\item Interpret the result in part (b) geometryically in terms of the tangent vector to a curve lying on a sphere in $\bb{R}^n$ centred at the origin.
	\end{enumerate}
\end{problem}

\pagebreak

\begin{problem}[27]
	Suppose $u, v, w \in V$. Prove that
	\[
		\norm{w - \frac{1}{2}(u + v)}^2 = \frac{\norm{w - u}^2 + \norm{w - v}^2}{2} - \frac{\norm{u - v}^2}{4}.
	\]
\end{problem}

\pagebreak

\begin{problem}[28]
Suppose $C$ is a subset of $V$ with the property that $u, v \in C$ implies $\frac{1}{2} (u + v) \in C$. Let $w \in V$. Show that there is at most one point in $C$ that is closest to $w$. In other words, show that there is at most one $u \in C$ such that
	\[
		\norm{w - u} \le \norm{w - v} \text{ for all } v \in C.
	\]
\end{problem}

\pagebreak

\begin{problem}[29]
    For \( u , v \in V \), define \( d(u, v) = \Vert u - v \Vert \)
    \begin{enumerate}
        \item Show that $d$ is a metric on $V$.
        \item Show that if $V$ is finite-dimensional, then $d$ is a complete metric on $V$ (meaning that every Cauchy sequence converges).
        \item Show that every finite-dimensional subspace of $V$ is a closed subset of $V$.
    \end{enumerate}
\end{problem}

\begin{solution}
    \begin{enumerate}
        \item In order to show that $d$ is a metric, we need to show that it is positive definite, symmetric, and satisfies the triangle inequality.
		
        First we show that $d$ is positive definite. Let \( u, v, w \) be any vectors in $V$. Since \( \Vert w \Vert = 0 \) if and only if \( w = 0 \), we have that \( d (u, v) = 0 \) if and only if \( u - v = 0 \). Thus, we have \( d (u, v) = 0 \) if and only if \( u = v \). Note also that \( \Vert w \Vert \ge 0 \), giving us \( d(u, v) \ge 0 \) always. Thus, $d$ is positive definite. 
	
        Now we show that $d$ is symmetric. This follows obviously from the properties of the norm of a scaled vector. For any \( u, v \in V \)
        \[
        \begin{split}
		d(u, v) & = \Vert u - v \Vert \\
			& = \Vert (-1) (v - u) \Vert \\
			& = \vert -1 \vert \Vert v - u \Vert \\
			& = \Vert v - u \Vert \\
			& = d(v, u).
        \end{split}
        \]\\
        Finally, the triangle inequalities for norms and metrics are equivalent. For any \( u, v, w \in V \) 
        \[
        d(u, v) + d(v, w) = \Vert v - u \Vert + \Vert w - v \Vert.
        \]
        By the triangle inequality (for norms) we have
        \[
        	\begin{split}
			\norm{v - u} + \norm{w - v} & \le \norm{(v - u) + (w - v)} \\
						    & = \norm{w - u}.
        	\end{split}
        \]
        Thus, we have the triangle inequality for metrics.
        \[
        d(u, v) + d(v, w) \le d(u, w).
        \]
        \item If $V$ is finite-dimensional, then we have an orthonormal basis \(e_1, \ldots, e_m\). Let $\set{v_n}$ be any Cauchy sequence in $V$.
		
        Then we can write each term of the sequence as
        \[
        v_n = a_{n, 1} e_1 + \ldots + a_{n, m} e_m.
        \]
	
        Since $\set{v_n}$ is Cauchy, we have that for any \( \varepsilon > 0 \), there exists some \( N \in \bb{N} \) such that for all \( j, k > N \) we have \( d(v_j, v_k) < \varepsilon \). We can rewrite the expression as
        \[
        	d(a_{j, 1} e_1 + \ldots + a_{j, m} e_m, a_{k, 1} e_1 + \ldots + a_{k, m} e_m) = \norm{ (a_{j, 1} - a_{k, 1}) e_1 + \ldots + (a_{j, m} - a_{k, m}) e_m}.
        \]
        Thus,
        \[
	        \norm{ (a_{j, 1} - a_{k, 1}) e_1 + \ldots + (a_{j, m} - a_{k, m}) e_m} < \varepsilon
        \]
	which reduces (by Pythagoras' theorem) to
	\[
        	\abs{a_{j, 1} - a_{k, 1}}^2 + \ldots + \abs{a_{j, m} - a_{k, m}}^2 < \varepsilon^2
        \]
	(with all the same qualifiers for \( j, k, \varepsilon \) ). 

	Since all the terms on the left hand side are non-negative, we can write
        \[
        \abs{a_{j, p} - a{k, p}}^2 < \varepsilon^2
        \]
        for any \( p \in \set{1, \ldots, m} \).
        Again, since both sides of the inequality are positive, we get
        \[
        \abs{a_{j, p} - a_{k, p}} < \varepsilon.
        \]
        Specifically, for any \( p \in \set{1, \ldots, m} \), for any \( \varepsilon > 0 \), there exists some \( N \in \bb{N} \) such that for all \( j, k > N \)
        \[
        \abs{a_{j, p} - a_{k, p}} < \varepsilon.
        \]
        That is, for any \( p \in \set{1, \ldots m} \) the sequence \( \set{a_{n, p}} \) is Cauchy.

        Since $V$ is a vector space over $\bb{R}$ or $\bb{C}$, both of which are complete, for all $p$, $\set{a_{n, p}}$ converges to some $A_p$. That is, for any \( \varepsilon' > 0 \), there is some \( N \in \bb{N} \) such that for all \( n > N \), \( \abs{A_p - a_{n, p}} < \varepsilon' \). We will show that the sequence of vectors converges to \( A_1 e_1 + \ldots + A_m e_m \). 

        For any \( \varepsilon' > 0 \) we can have some $N$ such that for all $n > N$, for all $p$, \( \abs{A_p - a_{n, p}} < \frac{\varepsilon}{m}' \). Specifically, if for each sequence of coefficients we get $N_p$, then $N = \max N_{p}$. We use Pythagoras' theorem again to get
        \[
        \begin{split}
		d(A_1 e_1 + \ldots + A_m e_m, a_{n, 1} e_1 + \ldots a_{n, m} e_m) & \le \abs{A_1 - a_{n, 1}} + \ldots + \abs{A_m - a_{n, m}} \\
										  & < m \frac{\varepsilon'}{m} = \varepsilon'
        \end{split}        \]
        Specifically, for any \( \varepsilon' > 0 \) there is some \( N \in \bb{N} \) such that for all \( n > N \)
        \[
        d(A_1 e_1 + \ldots + A_m e_m, a_{n, 1} e_1 + \ldots a_{n, m} e_m) < \varepsilon'.
        \]
        That is, \( \set{v_n} \) converges to \( A_1 e_1 + \ldots A_m e_m \).
        \item By the construction in part (b), any Cauchy sequence in a subspace $U$ of $V$ converges to a point in $U$. Thus, $U$ contains all of its limits points, and thus, is closed.
    \end{enumerate}
\end{solution}

\pagebreak

\begin{problem}[30]
	Fix a positive integer $n$. The \emph{Laplacian} $\Delta p$ of a twice differentiable function $p$ on $\bb{R}^n$ is the function on $\bb{R}^n$ defined by
	\[
		\Delta p = \frac{\partial^2 p}{\partial x_1} + \ldots + \frac{\partial^2 p}{\partial x_n}
	\]
A function $p$ is called \emph{harmonic} if \( \Delta p = 0 \).

A \emph{polynomial} on $\bb{R}^n$ is a linear combination of functions of the form \( x_1^{m_1} \cdots x_n^{m_n} \), where \( m_1 \ldots m_n \) are non-negative integers.

	Suppose $q$ is a polynomial on $\bb{R}^n$. Prove that there exists a harmonic polynomial $p$ on $\bb{R}^n$ such that \( p(x) = q(x) \) for every \( x \in \bb{R}^n \) with \( \Vert x \Vert = 1 \).
\end{problem}

\begin{solution}
	Let $\mathcal{P}(\bb{R}^n)$ denote the space of polynomials on $\bb{R}^n$ and let $\mathcal{P}_m(\bb{R}^n)$ be the space of polynomials on $\bb{R}^n$ of degree at most $m$.
	If $q \in \mathcal{P}(\bb{R}^n)$, and we want some harmonic polynomial $p$ such that \( p = q \) for every \( x \in \bb{R}^n \) with \( \norm{x} = 1 \), we can reasonably think that
	\[
		p = q + (1 - \norm{x}^2)r
	\]
	where $r$ is some polynomial. This is because $1 - \lVert x \rVert^{2}$ is effectively a polynomial anti-indicator function for the sphere --- zero everywhere on the sphere, and non-zero everywhere else. 

	Such a $p$ would necessarily give us \( p = q \) at all the desired points, but there might not be such a harmonic polynomial. Thus, we need to show that there is some harmonic polynomial of the same form. That is, we need to show that there exists some $r$ such that
	\[
		\begin{split}
			\Delta p & = \Delta q + \Delta ((1 - \norm{x}^2)r) \\
				 & = 0.
		\end{split}
	\]
	This amounts to showing that there exists some $r$ such that
	\[
		\Delta ((1 - \norm{x}^2)r) = -\Delta q.
	\]
	Notice that \( T : r \mapsto \Delta ((1 - \norm{x}^2)r) \) is a linear map that preserves the degree of the polynomial (multiplying by $\lVert x \rVert^{2}$ increases the degree by $2$, $\Delta$ decreases the degree by $2$, and both are linear). Thus, we can think of $T$ as a linear operator in $\mathcal{L}(\mathcal{P}_{m}(\mathbb{R}^{n}))$ and it will suffice to show that it is injective. 
		
	Suppose \( Ts = 0 \) for some polynomial $s$. Then
	\[
		\Delta ((1 - \norm{x}^2) s) = 0
	\]
	That is, $(1 - \lVert x \rVert^{2}) s$ is harmonic. Notice that $(1 - \lVert x \rVert^{2}) s = 0$ on the sphere as well. Thus, it must be $0$ everywhere (this is a fact about harmonic functions). Since $1 - \lVert x \rVert^{2}$ is non-zero everywhere that is not the sphere, $s = 0$ everywhere that is not the sphere, and thus, by its analyticity, is the zero polynomial. That is, $s = 0$. 
	
	This gives us that $T$ is injective and thus, surjective. Then, there exists some $r$ with $Tr = - \Delta q$, giving us the $p$ we want as
	\[
		p = (1 - \norm{x}^2) r + q
	\]
	for that $r$.
\end{solution}

\pagebreak

\begin{problem}[31]
	Use inner products to prove Apollonius's identity: in a triangle with sides of length $a$, $b$, and $c$, let $x$ be the length of the line segment from the midpoint of the side of length $c$ to the opposite vertex. Then 
	\[
		a^2 + b^2 = \frac{1}{2} c^2 + d^2
	\]
\end{problem}

\begin{solution}
	Consider the obvious choice of $v_a$, $v_b$, $v_c$, and $v_d$ such that \( \norm{v_\alpha} = \alpha \). Note that this then gives us
	\[
		\frac{v_c}{2} + v_d = v_a = v_c + v_b
	\]
	and
	\[
		\frac{v_c}{2} + v_b = v_d.
	\]
	
	Thus, we have 
	\[
		\norm{\frac{v_c}{2} + v_d}^2 = \norm{v_a}^2
	\]
	which we can expand using the dot product to get
	\[
		\frac{c^2}{4} + d^2 + \frac{v_c}{2} \cdot v_d = a^2.
	\]
	We also have 
	\[
		\norm{\frac{v_c}{2} + v_b}^2 = \norm{v_d}^2
	\]
	which gives us
	\[
		\frac{c^2}{4} + b^2 + \frac{v_c}{2} \cdot v_b = d^2.
	\]
	
	Swapping sides and adding the equations gives us
	\[
		2d^2 + \frac{v_c}{2} \cdot v_d = a^2 + b^2 + \frac{v_c}{2} \cdot v_b
	\]
	which we can rewrite as 
	\[
		2d^2 + \frac{v_c}{2} \cdot (v_d - v_b) = a^2 + b^2
	\]
	\[
		2d^2 + \frac{v_c}{2} \cdot v_c = a^2 + b^2
	\]
	which finally gives us
	\[
		a^2 + b^2 = 2d^2 + \frac{c^2}{2}
	\]
	as desired.
\end{solution}

% Add pairs of problems and solutions as needed

\end{document}
